\documentclass[a4paper,12pt]{article}
\usepackage[T2A]{fontenc}
\usepackage[utf8]{inputenc}
\usepackage[russian]{babel}
\usepackage{amsmath,amssymb,mathtools}
\usepackage{helvet}
\renewcommand{\familydefault}{\sfdefault}
\usepackage{icomma}
\usepackage[a4paper,margin=2.1cm]{geometry}
\usepackage{microtype}
\usepackage{tikz}
\usetikzlibrary{calc,positioning}
\usepackage{tabularx,booktabs}
\usepackage{enumitem}
\usepackage{hyperref}
\usepackage{parskip}
\usepackage{fancyhdr}
\usepackage{setspace}
\usepackage{xcolor}

\definecolor{accent}{HTML}{008080}
\definecolor{darkaccent}{HTML}{004D4D}
\definecolor{textgray}{HTML}{333333}
\definecolor{softgray}{HTML}{F2F4F4}

\hypersetup{colorlinks=true,linkcolor=darkaccent,urlcolor=darkaccent}
\onehalfspacing
\setlength{\parindent}{0pt}
\setlist[itemize]{leftmargin=1.7em,itemsep=0.25em,topsep=0.3em}
\setlist[enumerate]{leftmargin=2em,itemsep=0.55em,topsep=0.4em}

\pagestyle{fancy}
\fancyhf{}
\fancyhead[L]{\small\textcolor{textgray}{Теория вероятностей}}
\fancyhead[R]{\small\textcolor{textgray}{Ксения Клыкова}}
\fancyfoot[C]{\small\thepage}
\renewcommand{\headrulewidth}{0.4pt}
\newcommand{\sectionrule}{\vspace{-0.4em}\textcolor{accent}{\rule{\linewidth}{0.7pt}}\vspace{0.4em}}
\newcommand{\key}[1]{\textcolor{darkaccent}{\textbf{#1}}}
\newcommand{\prob}[1]{\operatorname{P}\!\left(#1\right)}

\begin{document}

\begin{titlepage}
\centering
\vspace*{2.1cm}
{\Large\textcolor{darkaccent}{\textbf{Чек-лист}}\par}
\vspace{0.7cm}
{\Huge\textbf{Дерево вероятностей}\par}
\vspace{0.35cm}
{\Large\textcolor{textgray}{Маршруты, условные вероятности и диагностические тесты}\par}
\vspace{2.2cm}
{\large Лёвин Артём Александрович\par}
\vspace{0.2cm}
{\large эксклюзивно для Ксении Клыковой\par}
\vfill
{\large 7 сентября 2026\par}
\vspace{1.7cm}
\begin{tikzpicture}[remember picture,overlay]
\draw[accent,line width=1.2pt]
($(current page.south west)+(1.7cm,1.8cm)$)
.. controls ($(current page.south west)+(5.4cm,2.8cm)$)
and ($(current page.south east)+(-5.4cm,0.8cm)$)
.. ($(current page.south east)+(-1.7cm,1.8cm)$);
\end{tikzpicture}
\end{titlepage}

\section*{1. Что показывает дерево вероятностей}
\sectionrule
Дерево вероятностей удобно использовать, когда случайный процесс проходит \key{несколько последовательных этапов}. Каждый узел показывает состояние после очередного этапа, а каждая ветвь --- один возможный переход.
\begin{itemize}
\item Из одного узла должны выходить все возможные продолжения процесса.
\item Сумма вероятностей ветвей, выходящих из одного узла, равна \(1\).
\item Полный путь от начала дерева до конечного состояния называют \key{траекторией}.
\end{itemize}

\section*{2. По пути умножаем, подходящие пути складываем}
\sectionrule
Если требуется, чтобы последовательно произошли события \(A\), затем \(B\), затем \(C\), то
\[
\prob{A\cap B\cap C}=\prob{A}\,\prob{B\mid A}\,\prob{C\mid A\cap B}.
\]
На дереве это означает: \key{вероятности ветвей одного маршрута перемножаются}.

Если нужный результат достигается по взаимно исключающим траекториям \(T_1,T_2,\ldots,T_k\), то
\[
\prob{T_1\cup\cdots\cup T_k}=\prob{T_1}+\cdots+\prob{T_k}.
\]
То есть \key{вероятности подходящих несовместимых маршрутов складываются}.

\subsection*{Мини-пример: погода}
Сегодня погода замечательная. Вероятность того, что на следующий день она останется такой же, равна \(0{,}8\); вероятность смены состояния равна \(0{,}2\). Чтобы через два дня погода снова была замечательной, подходят пути
\[
\text{З}\to\text{З}\to\text{З}
\quad\text{или}\quad
\text{З}\to\text{Х}\to\text{З}.
\]
Поэтому
\[
\prob{\text{З через два дня}}=0{,}8\cdot0{,}8+0{,}2\cdot0{,}2=0{,}68.
\]

\section*{3. Как читать длинное условие}
\sectionrule
\begin{enumerate}
\item Определите первый этап процесса.
\item Для каждого узла перечислите все возможные продолжения.
\item Подпишите вероятности ветвей.
\item Проверьте: сумма вероятностей из каждого узла равна \(1\).
\item Отметьте конечные состояния, удовлетворяющие вопросу.
\item Вдоль каждого подходящего пути перемножьте вероятности.
\item Для нескольких подходящих несовместимых путей сложите полученные произведения.
\end{enumerate}
В задачах по дням полезно подписывать даты или номера шагов рядом с уровнями дерева.

\section*{4. Повторный анализ после отрицательного результата}
\sectionrule
Пусть для больного пациента вероятность ложноотрицательного результата анализа равна \(0{,}3\). Тогда вероятность положительного результата одного анализа равна \(0{,}7\). Если после отрицательного результата анализ проводят ещё раз при тех же вероятностях, то
\[
\prob{\text{выявить за два анализа}}=0{,}7+0{,}3\cdot0{,}7=0{,}91.
\]
Или через противоположное событие:
\[
\prob{\text{выявить за два анализа}}=1-0{,}3^2=0{,}91.
\]

\newpage
\section*{5. Истинное состояние и результат теста}
\sectionrule
Обозначим
\[
D=\{\text{пациент действительно болен}\},\qquad +=\{\text{анализ положительный}\}.
\]
Тогда
\begin{itemize}
\item \(\prob{+\mid D}\) --- чувствительность теста;
\item \(\prob{+\mid\overline D}\) --- вероятность ложноположительного результата;
\item \(\prob{-\mid D}\) --- вероятность ложноотрицательного результата;
\item \(\prob{D}\) --- доля действительно больных в рассматриваемой группе.
\end{itemize}
\[
\boxed{\text{ложноположительный} = \text{здоров, но тест }+}
\]
\[
\boxed{\text{ложноотрицательный} = \text{болен, но тест }-}
\]

\section*{6. Вероятность положительного результата}
\sectionrule
Положительный результат может появиться по двум несовместимым путям:
\[
D\to +\qquad\text{или}\qquad\overline D\to +.
\]
Следовательно,
\[
\boxed{\prob{+}=\prob{D}\prob{+\mid D}+\prob{\overline D}\prob{+\mid\overline D}}.
\]
Если \(\prob{D}=0{,}05\), \(\prob{+\mid D}=0{,}9\), \(\prob{+\mid\overline D}=0{,}01\), то
\[
\prob{+}=0{,}05\cdot0{,}9+0{,}95\cdot0{,}01=0{,}0545.
\]

\section*{7. Обратная задача: найти долю больных}
\sectionrule
Пусть \(\prob{D}=x\), \(\prob{+\mid D}=a\), \(\prob{+\mid\overline D}=b\), \(\prob{+}=r\). Тогда
\[
ax+b(1-x)=r,
\]
откуда при \(a\ne b\)
\[
\boxed{x=\frac{r-b}{a-b}}.
\]
Например, при \(r=0{,}06\), \(a=0{,}9\), \(b=0{,}01\):
\[
x=\frac{0{,}05}{0{,}89}\approx0{,}0562.
\]

\section*{8. Вероятность болезни после положительного теста}
\sectionrule
Если тест уже положительный, требуется условная вероятность
\[
\prob{D\mid +}=\frac{\prob{D}\prob{+\mid D}}{\prob{+}}.
\]
Для примера с \(\prob{D}=0{,}05\), чувствительностью \(0{,}9\) и ложноположительной вероятностью \(0{,}01\):
\[
\prob{D\mid +}=\frac{0{,}05\cdot0{,}9}{0{,}0545}\approx0{,}8257.
\]
Важно различать \(\prob{D}\) --- долю больных во всей группе --- и \(\prob{D\mid +}\) --- вероятность болезни среди людей с положительным тестом.

\section*{9. Быстрый алгоритм}
\sectionrule
\begin{enumerate}
\item Что является первым этапом процесса?
\item Какие исходы возможны на первом этапе?
\item Какие исходы возможны после каждого из них?
\item Сумма вероятностей из каждого узла равна \(1\)?
\item Какие конечные ветви удовлетворяют вопросу?
\item Вдоль каждого подходящего пути перемножьте вероятности.
\item Несколько подходящих несовместимых путей сложите.
\item Проверьте, какую именно вероятность спрашивают: \(\prob{+}\), \(\prob{D}\) или \(\prob{D\mid +}\).
\end{enumerate}

\section*{10. Типичные ошибки}
\sectionrule
\begin{itemize}
\item Лишний уровень дерева: подписывайте даты или номера этапов.
\item Перепутаны умножение вдоль пути и сложение разных подходящих путей.
\item Перепутаны истинное состояние и результат теста.
\item Неверно прочитан ложный результат.
\item Нарушена сумма вероятностей ветвей из одного узла.
\item Смешаны вопросы о доле больных и о вероятности болезни при положительном тесте.
\item Правильная модель испорчена арифметикой.
\end{itemize}

\newpage
\section*{Тренировочный блок}
\sectionrule
\begin{enumerate}
\item Сегодня погода ясная. Вероятность того, что на следующий день погода сохранится, равна \(0{,}75\), а изменится --- \(0{,}25\). Найдите вероятность того, что через два дня погода снова будет ясной.
\item У больного пациента анализ даёт отрицательный результат с вероятностью \(0{,}2\). После отрицательного результата анализ повторяют один раз при тех же вероятностях. Найдите вероятность того, что заболевание будет выявлено хотя бы одним из двух анализов.
\item В группе пациентов \(4\%\) действительно больны. Чувствительность теста равна \(0{,}92\), вероятность ложноположительного результата --- \(0{,}03\). Найдите вероятность положительного результата у случайно выбранного пациента.
\end{enumerate}

\newpage
\section*{Ответы к тренировочному блоку}
\sectionrule
\renewcommand{\arraystretch}{1.35}
\begin{table}[h!]
\centering
\caption{Краткие ответы}
\begin{tabularx}{\linewidth}{@{}cX@{}}
\toprule
\textbf{№} & \textbf{Ответ}\\
\midrule
1 & \(0{,}625\)\\
2 & \(0{,}96\)\\
3 & \(0{,}0656\)\\
\bottomrule
\end{tabularx}
\end{table}

\end{document}
